\pagebreak

\section{CR Alternativi:
CyberPatriot}\label{cr-alternativi-cyberpatriot}

È un programma di educazione nazionale per il K-12 creato dalla
\emph{Air \& Space Forces Association} (AFA) \footnote{\label{fn:ref_12}Air \& Space
  Forces Association, ``CyberPatriot: The National Youth Cyber Education
  Program'', https://www.uscyberpatriot.org}. I team assumono il ruolo
di professionisti di cybersecurity che hanno l'obiettivo di proteggere
vari sistemi entro un limite di tempo, analizzando le macchine virtuali,
identificandone le vulnerabilità e mitigandole.

Ciascun team riceve una o più macchine virtuali complete pre-configurate
(basate su distribuzioni Windows Server/Desktop o Linux), volutamente
affette da molteplici difetti di sicurezza: account non autorizzati,
policy di password permissive, servizi di rete non necessari in
esecuzione, permessi errati sul filesystem e malware dormiente. I
partecipanti operano come amministratori di sistema (``Blue Team''), con
il compito di individuare e correggere tali falle entro una finestra
temporale prefissata, mantenendo al contempo operativi i servizi critici
per il business.

Il CyberPatriot Scoring System (CSS) tiene traccia della progressione
del team e dei punteggi dei team. Un demone locale esegue controlli
periodici sulla macchina per determinarne lo stato, mentre un server
centrale registra i punteggi comunicando con i demoni delle macchine.

Questo sistema tuttavia è \textbf{closed-source} e solo l'\textbf{AFA}
può utilizzarlo.

\begin{longtable}[]{@{}
  >{\raggedright\arraybackslash}p{(\linewidth - 4\tabcolsep) * \real{0.2212}}
  >{\raggedright\arraybackslash}p{(\linewidth - 4\tabcolsep) * \real{0.3540}}
  >{\raggedright\arraybackslash}p{(\linewidth - 4\tabcolsep) * \real{0.4248}}@{}}
\toprule\noalign{}
\begin{minipage}[b]{\linewidth}\raggedright
Caratteristica
\end{minipage} & \begin{minipage}[b]{\linewidth}\raggedright
CyberPatriot
\end{minipage} & \begin{minipage}[b]{\linewidth}\raggedright
BlueAgent
\end{minipage} \\
\midrule\noalign{}
\endhead
\bottomrule\noalign{}
\endlastfoot
\textbf{Virtualizzazione} & Macchine Virtuali (VMware / VirtualBox) &
Container di sistema (Incus) \\
\textbf{Consumo di Risorse} & Diversi GB di RAM per singola VM &
\textasciitilde50 MB di RAM per container \\
\textbf{Motore di Valutazione} & Demone locale (CSS) all'interno della
VM & Oracolo esterno in sandbox di rete isolata (OVN) \\
\end{longtable}
